% ── §7g Methods: Comprehensive Method Comparison ──
\section{Comprehensive Method Comparison}
\label{sec:methods-comparison}

\subsection{Methods Evaluated}
We evaluate seven distinct methods for contact prediction on the PSICOV150
benchmark, spanning three methodological families:

\begin{enumerate}[nosep]
	\item \textbf{Karnaugh-map circuits}: Level-1 (group-level) and Level-2
	      (identity-level) Quine--McCluskey-minimised circuits with cell-rate
	      scoring.
	\item \textbf{Continuous statistics}: plain mutual information,
	      APC-corrected mutual information (MIp), perplexity ratio, and
	      APC-corrected Frobenius-norm DCA from our mean-field implementation.
	\item \textbf{Literature baseline}: published contact predictions from the
	      PSICOV method (sparse inverse covariance estimation), downloaded
	      directly from the benchmark supplementary data.
\end{enumerate}

\subsection{Results Across Separation Bands}
Contact prediction performance varies substantially across separation bands.
We report precision@L/5 stratified by CASP convention bands:

\begin{table}[ht]
	\centering
	\caption{Precision@L/5 stratified by CASP separation band (mean over held-out proteins).}
	\label{tab:bands}
	\begin{tabular}{lccc}
		\toprule
		Method           & Short ($6$--$11$) & Medium ($12$--$23$) & Long ($\geq 24$) \\
		\midrule
		Circuit-L2       & .179              & .134                & .137             \\
		Circuit-L1       & .135              & .110                & .116             \\
		Plain MI         & .115              & .081                & .077             \\
		Perplexity ratio & .052              & .036                & .031             \\
		Separation-only  & .033              & .052                & .051             \\
		PSICOV published & .409              & .478                & .637             \\
		\bottomrule
	\end{tabular}
\end{table}

The circuit maintains its advantage over MI in every band, though the margin
narrows at long range.  The PSICOV baseline shows the opposite trend,
achieving its highest precision at long range where covariation methods are
typically evaluated; this reflects the sparse inverse covariance estimator's
ability to remove indirect correlations that are most prevalent at short to
medium range.

\subsection{Computational Cost}
All experiments were performed on a shared server with an NVIDIA A100 GPU
(80~GB HBM2e) and 24-core Intel Xeon CPU.  The per-protein wall-clock times
are:

\begin{center}
	\begin{tabular}{lr}
		\toprule
		Operation                             & Time           \\
		\midrule
		Feature extraction (per protein)      & $< 1$~s        \\
		MI matrix, full length (GPU)          & $0.1$~s        \\
		mfDCA covariance inversion (CPU)      & $5$--$60$~s    \\
		QM minimisation (1024 cells)          & $<$ 1~s        \\
		Full PATH-B evaluation (150 proteins) & $\approx$ 21~s \\
		\bottomrule
	\end{tabular}
\end{center}
